
\documentclass[12pt]{article}
\usepackage[margin=1in]{geometry}
\usepackage{setspace}
\usepackage{amsmath}
\usepackage{natbib}

\title{Political Polarization and the Interpretation of Auditor Signals}
\author{Neil}
\date{}

\begin{document}
\maketitle
\onehalfspacing

\section{Theory}

We model investors receiving a public audit signal $s \in \{0,1\}$ regarding firm quality. Investors have prior beliefs $\pi_i$ that differ across ideological groups. Following \citet{kandel1995}, investors interpret signals through heterogeneous priors:

\begin{equation}
P_i(\theta=1|s) = \frac{P(s|\theta=1)\pi_i}{P(s)}
\end{equation}

Political polarization increases dispersion in $\pi_i$, leading to greater disagreement in posterior beliefs.

This generates two key predictions:

\begin{itemize}
\item Greater trading volume due to belief heterogeneity
\item Larger absolute price movements due to disagreement
\end{itemize}

When audit signals are more ambiguous, the likelihood function $P(s|\theta)$ is less precise, amplifying the role of priors and thus polarization effects.

\section{Implications}

Polarization does not change the information content of audit signals per se, but changes how that information is processed. This implies that identical audit disclosures can lead to different market outcomes depending on the political composition of investors.

\bibliographystyle{apalike}
\bibliography{references}

\end{document}
